\documentclass[10pt,a4paper]{amsart}
\usepackage[margin=19mm]{geometry}
\usepackage{amsmath,amssymb,mathtools}
\usepackage[scr=rsfs,cal=euler]{mathalfa}
\usepackage{xfrac}
\usepackage[unicode,hidelinks,bookmarks=false]{hyperref}
\newcommand{\ie}{i.e.\ }
\newcommand{\cf}{cf.\ }
\newcommand{\resp}{respectively\ }
\newcommand{\Subsection}{\subsection}
\providecommand{\nobreakdash}{\nobreak\hbox{-}}
\setlength{\emergencystretch}{2em}
\multlinegap0pt
\usepackage{bm}
\usepackage{bbm}
\usepackage{dsfont}
\usepackage{xspace}
\usepackage{cancel}
\usepackage{mathtools}
\usepackage{cleveref}
\usepackage{amsthm}
\makeatletter
\@ifundefined{thm}{\newtheorem{thm}{Theorem}}{}
\@ifundefined{cor}{\newtheorem{cor}[thm]{Corollary}}{}
\@ifundefined{lemma}{\newtheorem{lemma}[thm]{Lemma}}{}
\@ifundefined{prop}{\newtheorem{prop}[thm]{Proposition}}{}
\makeatother
\newcommand{\C}{\mathbb{C}}
\newcommand{\N}{\mathbb{N}}
\newcommand{\Z}{\mathbb{Z}}
\title[Twisted Knots and the Perturbed Alexander Invariant]{Twisted Knots and the Perturbed Alexander Invariant}
\author{Joe Boninger}
\date{Source: arXiv:2403.03754v4 (CC BY 4.0)}
\begin{document}
\maketitle

\noindent\emph{Source and licence.} Twisted Knots and the Perturbed Alexander Invariant. Version arXiv:2403.03754v4 (CC BY 4.0), distributed under the Creative Commons Attribution 4.0 International licence, as recorded in the arXiv licence metadata for this exact version and shown on the abstract page https://arxiv.org/abs/2403.03754v4. Full rights notice: licenses/arxiv-2403.03754-CC-BY-4.0.txt.\par\smallskip

\noindent\textit{Editorial scope.} Counted unit: the Thm whose source label is \texttt{prop:main} in the article (source0.tex lines 1186--1197, source label \texttt{prop:main}), delivered together with the article's own complete original proof of it (source0.tex lines 1199--1309, one \texttt{proof} environment of 111 lines). No mathematics is written, repaired or re-derived: statement and proof are the article's own text line for line. Required context retained uncounted: thm:cmc (lines 444--454); cor:g\_matrix (lines 507--514); thm:cmc\_det (lines 518--525); lem:stab\_cont (lines 741--751); lem:alex (lines 806--809); eq:newphi (lines 1049--1052); eq:new\_rho (lines 1054--1057); lem:tech (lines 1088--1095). Every definition, display and label of those lines is retained unchanged. Omitted from this package and not redistributed: 1--443, 455--506, 515--517, 526--740, 752--805, 810--1048, 1053--1053, 1058--1087, 1096--1185, 1198--1198, 1310--1685. Nothing inside a retained excerpt is omitted or altered: every retained body is delivered exactly as the article's lines are written, so no omission receipt of any kind accompanies this package. Citations: the retained excerpts carry no citation command at all, so no bibliography entry of the article is transcribed and no reference is redistributed. Reference adaptation: none was needed and none was made; every reference inside the delivered unit resolves inside it, so no unresolved reference or citation is produced. Printed slips kept literal: no printed word, symbol, hypothesis or number of the counted statement or of its proof was altered. Layout and class: the article's own class is \texttt{nyjm} with packages including amsfonts, amstext, amsmath, amsthm, amscd, amssymb, graphicx, color, cite, hyperref, enumitem, setspace, float, subcaption; the wrapper is the lane's pinned \texttt{amsart} editorial wrapper at 10pt on A4, which restates the theorem-like environments the retained text needs and adds the article's own macro definitions that the retained text needs (\texttt{\textbackslash C}, \texttt{\textbackslash N}, \texttt{\textbackslash Z}), with extra wrapper packages (bm, bbm, dsfont, xspace, cancel, mathtools, cleveref). The delivered numbering is the unit's own. Assets: no retained span contains a figure environment, a graphics-inclusion command, a PDF-page inclusion, a table or a TikZ picture; no image file is copied into this package, the package declares no asset dependency and no image file is redistributed with it. Licence: version arXiv:2403.03754v4 of this article is distributed under the Creative Commons Attribution 4.0 International licence, as recorded in the arXiv licence metadata for this exact version and shown on the abstract page https://arxiv.org/abs/2403.03754v4. A full rights notice is delivered at licenses/arxiv-2403.03754-CC-BY-4.0.txt. Characters: every retained excerpt is pure ASCII, so the delivered engine, XeLaTeX with its default Latin Modern text font, reports no missing character.
	\begin{prop}
		\label{thm:cmc}
		The Green's function of $D/U$,
		$$
		g_{D/U} : \mathcal{S}(D/U) \times \mathcal{S}(D/U) \to \Z(T),
		$$
		 is well-defined for all $T \in \C$ in an open neighborhood of $1$. For all $s,t \in \mathcal{S}(D/U) \subset \mathcal{S}(D)$,
		 $$
		 	g_{D/U}(s,t) = g_D(s,t).
		 $$
	\end{prop}

	\begin{cor}
		\label{cor:g_matrix}
		For any contraction as above, the Green's matrix
		$$
		G_{D/U} = (I - A_{D/U})^{-1} \in M_{|\mathcal{S}(D/U)|}(\Z(T))
		$$
		is well defined, and its entries coincide with the Green's function $g_{D/U}$ as in the case of tangle Markov chains.
	\end{cor}

	\begin{prop}
		\label{thm:cmc_det}
		If the contracted region $U$ admits no cycles, then
		$$
		\det(I - A_D) = \det(I - A_{D/U}).
		$$
		Equivalently, $\det(G_D) = \det(G_{D/U})$.
	\end{prop}

	\begin{lemma}
		\label{lem:stab_cont}
		If $\{p_t\}_{t \in \N}$ and $\{q_t\}_{t \in \N}$ are two sequences of Laurent polynomials which stabilize positively, then the sequences $\{p_t + q_t\}_t$ and $\{p_tq_t\}_t$ stabilize positively as well. Furthermore,
		$$
		\lim_{t \to \infty} (p_t + q_t) = \lim_{t \to \infty} p_t + \lim_{t \to \infty} q_t
		$$
		and
		$$
		\lim_{t \to \infty} (p_tq_t) = (\lim_{t \to \infty} p_t)(\lim_{t \to \infty} q_t).
		$$
	\end{lemma}

	\begin{lemma}
		\label{lem:alex}
		Let $\{K_t\}$ be a coherently oriented family of knots twisted along $n$ strands. Then the sequence $\{T^{tn(n -1)/2}\Delta_{K_t}(T)\}_t$ stabilizes positively to a rational function as $t \to \infty$.
	\end{lemma}

	\begin{equation}
		\label{eq:newphi}
		\tilde{\varphi}(s_k) = \tilde{\varphi}(k) =  \det(I - A)^2 \varphi_k (g_{kk} - 1/2) = \det(I - A) \varphi_k \tilde{g}_{kk} - \varphi_k\det(I - A)^2/2.
	\end{equation}

	\begin{equation}
		\label{eq:new_rho}
		\rho_1(K) = T^{-\varphi(D) - w(D)} \Big( \sum_{c \in \mathcal{C}} \tilde{R}_1(c) - \sum_{s_k \in \mathcal{S}} \tilde{\varphi}(k) \Big).
	\end{equation}

	\begin{lemma}
		\label{lem:tech}
		For all $r_0 \in \Z$ there exists $m > 0$ such that for any $t > 2m$, any full twist $\tau$ in $U_{t,m}^\text{mid} \subset D_t$ and any crossing $c_k(\tau)$ in $\tau$, we have
		$$
		\big(\widetilde{R}_1(c_k(\tau))\big)_r = \big(\widetilde{R}_1(c_k(\tau_\infty))\big)_r.
		$$
		for all $r \leq r_0$.
	\end{lemma}

	\begin{thm}
		\label{prop:main}
		Define
		$$
		d_t = d_t(\{K_t\}) = T^{tn(n - 1)} \Big( T^{n(n - 1)}\rho_1(K_{t + 1}) - \rho_1(K_t) \Big) \in \Z[T,T^{-1}]
		$$
		for all $t \in \N$. Then the sequence $\{d_t\}$ stabilizes positively with limit
		$$
		\lim_{t \to \infty} d_t = T^{-\varphi(D_0) -w(D_0)} \sum_{k = 1}^{n(n - 1)} \widetilde{R}_1(c_k(\tau_\infty)),
		$$
		where $\tau_\infty$ is the distinguished full twist in $D^\tau_\infty$.
	\end{thm}

	\begin{proof}
		Given an ends index $m$, for all $t > 2m$ we write $\mathcal{S}_{t,m}^\text{mid}$ to indicate $\mathcal{S}(U^\text{mid}_{t,m})$ and $\mathcal{S}_{t,m}^{\overline{\text{mid}}}$ to denote its complement $\mathcal{S}(D_t) - \mathcal{S}^\text{mid}_{t,m}$. Similarly, we write $\mathcal{C}_{t,m}^\text{mid}$ to mean the set of crossings $\mathcal{C}(U^\text{mid}_{t,m})$, and $\mathcal{C}_{t,m}^{\overline{\text{mid}}}$ for the complementary set $\mathcal{C}(D_t) - \mathcal{C}^\text{mid}_{t,m}$.
		
		We write
		$$
		\rho_1(K_t) = \rho^\text{mid}_m(K_t) + \rho^{\overline{\text{mid}}}_m(K_t),
		$$
		where as in (\ref{eq:new_rho})
		$$
		\rho^\text{mid}_m(K_t) = T^{-\varphi(D_t) -w(D_t)} \Big( \sum_{c \in \mathcal{C}_{t,m}^\text{mid}} \tilde{R}_1(c) - \sum_{k \in \mathcal{S}_{t,m}^\text{mid}} \tilde{\varphi}_k \Big),
		$$
		and $\rho^{\overline{\text{mid}}}_m$ is defined analogously for $\mathcal{C}_{t,m}^{\overline{\text{mid}}}$ and $\mathcal{S}_{t,m}^{\overline{\text{mid}}}$. Then
		\begin{equation}
			\label{eq:setup}
			d_t = T^{tn(n - 1)} \big(T^{n(n - 1)} \rho^\text{mid}_m(K_{t + 1}) -  \rho^\text{mid}_m(K_t)\big) + T^{tn(n - 1)}\big(T^{n(n - 1)} \rho^{\overline{\text{mid}}}_m(K_{t + 1}) -  \rho^{\overline{\text{mid}}}_m(K_t) \big).
		\end{equation}
		Fix $r_0 \in \Z$; we consider the left grouping in (\ref{eq:setup}) first. We observe that all strands in $U^\text{mid}_{t,m}$ for any $t$ and $m$ have zero turning number, so by (\ref{eq:newphi})
		$$
		\rho^\text{mid}_m(K_t) = T^{-\varphi(D_t) -w(D_t)} \sum_{c \in \mathcal{C}_{t,m}^\text{mid}} \tilde{R}_1(c).
		$$
		Additionally, since $\varphi(D_t) = \varphi(D_0)$ for all $t$ and
		$$
		w(D_t) = w(D_0) + tn(n - 1),
		$$
		we have
		\begin{equation}
			\label{eq:first_grp}
			T^{tn(n - 1)} \big(T^{n(n - 1)} \rho^\text{mid}_m(K_{t + 1}) -  \rho^\text{mid}_m(K_t)\big) = T^{-\varphi(D_0) -w(D_0)}\big(\sum_{c \in \mathcal{C}_{t + 1,m}^\text{mid}} \tilde{R}_1(c) - \sum_{c \in \mathcal{C}_{t,m}^\text{mid}} \tilde{R}_1(c)\big).
		\end{equation}
		Using Lemma \ref{lem:tech} we choose an ends index $m$ such that for any $t > 2m$, any full twist $\tau$ in $U^\text{mid}_{t,m}$ and any crossing $c_k(\tau)$,
		$$
		\big(\widetilde{R}_1(c_k(\tau))\big)_r = \big(\widetilde{R}_1(c_k(\tau_\infty))\big)_r
		$$
		for all $r \leq r_0$. This implies
		$$
		\big(\sum_{k = 1}^{n(n-1)} \widetilde{R}_1(c_k(\tau))\big)_r = \big(\sum_{k = 1}^{n(n-1)} \widetilde{R}_1(c_k(\tau_\infty))\big)_r
		$$
		for all $r\leq r_0$, so in fact
		$$
		\big( \sum_{c \in \mathcal{C}_{t,m}^\text{mid}} \tilde{R}_1(c) \big)_r = \big((t - 2m) \cdot \sum_{r = 1}^{n(n-1)} \widetilde{R}_1(c_r(\tau_\infty)) \big)_r
		$$
		for all $r \leq r_0$ and all $t > 2m$. From this and (\ref{eq:first_grp}), we conclude that
		\begin{equation}
			\label{eq:plugin}
			\Big( T^{tn(n - 1)} \big(T^{n(n - 1)} \rho^\text{mid}_m(K_{t + 1}) -  \rho^\text{mid}_m(K_t)\big) \Big)_r = \Big(T^{-\varphi(D_0) -w(D_0)} \sum_{k = 1}^{n(n-1)} \widetilde{R}_1(c_k(\tau_\infty)) \Big)_r
		\end{equation}
		for all $r \leq r_0$ and all $t > 2m$.
		
		We finish the proof of the theorem by showing the second grouping in (\ref{eq:setup}) stabilizes positively to zero. For this, define
		$$
		D^{\overline{\text{mid}}}_{t,m} = D_t/U^\text{mid}_{t,m}.
		$$
		Additionally, let $D^{\overline{\text{mid}}}_{\infty,m}$ be the diagram formed by replacing the region $U^\text{mid}_{t,m}$ with a single infinite twist vertex. We note that the choice of $t$ does not matter for this definition, so long as $t > 2m$. Let $A^{\overline{\text{mid}}}_{t,m}$, $A^{\overline{\text{mid}}}_{\infty,m}$, $G^{\overline{\text{mid}}}_{t,m}$ and $G^{\overline{\text{mid}}}_{\infty,m}$ be transition and Green's matrices for $D^{\overline{\text{mid}}}_{t,m}$ and $D^{\overline{\text{mid}}}_{\infty,m}$ respectively, and let
		\begin{align*}
		\widetilde{G}^{\overline{\text{mid}}}_{t,m} &= \text{adj}(I - A^{\overline{\text{mid}}}_{t,m}) \\
		\widetilde{G}^{\overline{\text{mid}}}_{\infty,m} &= \text{adj}(I - A^{\overline{\text{mid}}}_{\infty,m}).
		\end{align*}
	
		The diagrams $D_t - U^\text{mid}_{t,m}$ are identical for all $t > 2m$, and we identify the sets of states $\mathcal{S}_{t,m}^{\overline{\text{mid}}}$ with each other and with $\mathcal{S}(D^{\overline{\text{mid}}}_{\infty,m})$ in the obvious way. We choose an indexing for this set, and extend this to an indexing on $\mathcal{S}(D_t)$ for each $t$ via the inclusion $D_t - U^\text{mid}_{t,m} \hookrightarrow D_t$. Similarly, we identify the sets of crossings $\mathcal{C}^{\overline{\text{mid}}}_{t,m}$ and $\mathcal{C}(D^{\overline{\text{mid}}}_{\infty,m})$ for all $t$.
		
		As in the proof of Lemma \ref{lem:tech}, since the contracted region $U^\text{mid}_{t,m}$ admits no cycles, Propositons \ref{thm:cmc} and \ref{thm:cmc_det} and Corollary \ref{cor:g_matrix} give
		$$
		\det(I - A_t) = \det(I - A^{\overline{\text{mid}}}_{t,m})
		$$
		and
		$$
		(\widetilde{G}_t)_{ij} = \det(I - A_t)(G_t)_{ij} = \det(I - A^{\overline{\text{mid}}}_{t,m})(G^{\overline{\text{mid}}}_{t,m})_{ij} = (\widetilde{G}^{\overline{\text{mid}}}_{t,m})_{ij}
		$$
		for all $i$ and $j$ indexing states in $\mathcal{S}_{t,m}^{\overline{\text{mid}}}$. The same arguments used in the proofs of \Cref{lem:alex} and Lemma \ref{lem:tech} also show that
		\begin{equation}
			\label{eq:final_stab}
			\lim_{t \to \infty} A^{\overline{\text{mid}}}_{t,m} = A^{\overline{\text{mid}}}_{\infty,m},
		\end{equation}
		which implies 
		$$
		\lim_{t \to \infty} \widetilde{G}^{\overline{\text{mid}}}_{t,m} = \widetilde{G}^{\overline{\text{mid}}}_{\infty,m}
		$$
		by \Cref{lem:stab_cont}.
		
		Given a fixed crossing $c$ of $D_{2m + 1} - U^\text{mid}_{2m + 1,m}$, let $c^t$ denote the corresponding crossing of $\mathcal{C}^{\overline{\text{mid}}}_{t,m}$ and $c^\infty$ the same crossing in $\mathcal{C}(D^{\overline{\text{mid}}}_{\infty,m})$. By the above discussion, for any such $c$ with triple $(\sigma, i,j)$, we have
		\begin{align*}
		\lim_{t \to \infty} \widetilde{R}_1(c^t) &= \lim_{t \to \infty}\sigma((\widetilde{G}_t)_{ji}((\widetilde{G}_t)_{j^+, j} + (\widetilde{G}_t)_{j, j^+} \\
		& \ \ \ \ \ - (\widetilde{G}_t)_{ij}) - (\widetilde{G}_t)_{ii}((\widetilde{G}_t)_{j,j^+} - \det(I - A_t)) - \det(I - A_t)^2/2) \\
		&= \lim_{t \to \infty}\sigma((\widetilde{G}^{\overline{\text{mid}}}_{t,m})_{ji}((\widetilde{G}^{\overline{\text{mid}}}_{t,m})_{j^+, j} + (\widetilde{G}^{\overline{\text{mid}}}_{t,m})_{j, j^+} \\
		& \ \ \ \ \ - (\widetilde{G}^{\overline{\text{mid}}}_{t,m})_{ij}) - (\widetilde{G}^{\overline{\text{mid}}}_{t,m})_{ii}((\widetilde{G}^{\overline{\text{mid}}}_{t,m})_{j,j^+} - \det(I - A^{\overline{\text{mid}}}_{t,m})) - \det(I - A^{\overline{\text{mid}}}_{t,m})^2/2) \\
		&= \sigma((\widetilde{G}^{\overline{\text{mid}}}_{\infty,m})_{ji}((\widetilde{G}^{\overline{\text{mid}}}_{\infty,m})_{j^+, j} + (\widetilde{G}^{\overline{\text{mid}}}_{\infty,m})_{j, j^+} \\
		& \ \ \ \ \ - (\widetilde{G}^{\overline{\text{mid}}}_{\infty,m})_{ij}) - (\widetilde{G}^{\overline{\text{mid}}}_{\infty,m})_{ii}((\widetilde{G}^{\overline{\text{mid}}}_{\infty,m})_{j,j^+} - \det(I - A^{\overline{\text{mid}}}_{\infty,m})) - \det(I - A^{\overline{\text{mid}}}_{\infty,m})^2/2) \\
		&= \tilde{R}_1(c^\infty).
		\end{align*}
		Similarly, if $s^t$ is a fixed state of $D_t - U^\text{mid}_{t,m}$, let $s^\infty$ be the corresponding state of $\mathcal{S}(D^{\overline{\text{mid}}}_{\infty,m})$. Since the turning number $\varphi(s^t)$ does not depend on the value of $t$, a computation analogous to the one above shows
		$$
		\lim_{t \to \infty} \tilde{\varphi}(s^t) = \tilde{\varphi}(s^\infty).
		$$
		Returning to the second grouping of (\ref{eq:setup}), these two computations together imply that
		\begin{align*}
		&\lim_{t \to \infty} T^{n(n - 1)} \rho^{\overline{\text{mid}}}_m(K_{t + 1}) -  \rho^{\overline{\text{mid}}}_m(K_t) \\
		& \ \ \ \ \ = \lim_{t \to \infty}  T^{-w(D_0) - \varphi(D_0)}\big(\sum_{c \in \mathcal{C}_{{t+1},m}^{\overline{\text{mid}}}} \tilde{R}_1(c) - \sum_{k \in \mathcal{S}_{{t+1},m}^{\overline{\text{mid}}}} \tilde{\varphi}(k) - \sum_{c \in \mathcal{C}_{t,m}^{\overline{\text{mid}}}} \tilde{R}_1(c) + \sum_{k \in \mathcal{S}_{t,m}^{\overline{\text{mid}}}} \tilde{\varphi}(k) \big) \\
		& \ \ \ \ \ = T^{-w(D_0) - \varphi(D_0)}\big(\sum_{c \in \mathcal{C}(D^{\overline{\text{mid}}}_{\infty,m})} \tilde{R}_1(c) - \sum_{k \in \mathcal{S}(D^{\overline{\text{mid}}}_{\infty,m})} \tilde{\varphi}(k) - \sum_{c \in \mathcal{C}(D^{\overline{\text{mid}}}_{\infty,m})} \tilde{R}_1(c) + \sum_{k \in \mathcal{S}(D^{\overline{\text{mid}}}_{\infty,m})} \tilde{\varphi}(k) \big) \\
		& \ \ \ \ \ = 0.
		\end{align*}
		Thus we can choose $t_0 > 2m$ such that for all $t > t_0$ and $r \leq r_0$.
		$$
		\big( T^{n(n - 1)} \rho^{\overline{\text{mid}}}_m(K_{t + 1}) -  \rho^{\overline{\text{mid}}}_m(K_t) \big)_r = 0.
		$$
		
		Subsituting (\ref{eq:plugin}) and the above equation into (\ref{eq:setup}), we conclude that for all $t > t_0$ and all $r \leq r_0$,
		$$
		(d_t)_r = \Big(T^{-\varphi(D_0) -w(D_0)} \sum_{k = 1}^{n(n-1)} \widetilde{R}_1(c_k(\tau_\infty)) \Big)_r.
		$$
		Since $r_0$ was arbitrary, $d_t$ stabilizes positively to the desired limit.
	\end{proof}

\end{document}
